\documentclass[12pt]{article}
\usepackage[utf8]{inputenc}
\usepackage[T1]{fontenc}
\usepackage{hyperref}
\title{Tula Language Resources for K-12 Curriculum Development}
\author{Compiled by Claude Sonnet 5 (claude-sonnet-5) for Andr\'as Kornai}
\date{}
\begin{document}
\maketitle

\section{Introduction}
Tula (Kutule, Kitule; ISO 639-3 tul) is spoken in Gombe State, Nigeria, and belongs to the Tula--Waja group. The search found no textbooks, primers or school programmes for Tula. What exists is a small body of dictionaries, a historical grammar and academic descriptions. This report lists only what appeared in search results, and none of the items was opened or checked in full.

\section{Primary Educational Resources}
No primers, readers or school materials in Tula turned up in the search. The nearest thing to a teaching resource is the dictionary below.
\begin{itemize}
\item \emph{Modern Tula-English-Hausa Dictionary} by Mohammed Aminu Mu'azu and Hannatu Ibrahim Polo. A LINGUIST List notice (listed below under Linguistic and Reference Materials) describes it as the first dictionary of its kind, with English and Hausa glosses, example sentences and grammatical classes. Publisher and date were not confirmed. \url{https://linguistlist.org/issues/26/2013}
\end{itemize}

\section{Linguistic and Reference Materials}
\begin{itemize}
\item LINGUIST List 26.2013, the announcement of the Tula-English-Hausa dictionary. \url{https://linguistlist.org/issues/26/2013}
\item Grambank entry for Tula. It cites Blench and Tula (2012), \emph{A dictionary of Yi Kutul\textepsilon{} fi Swa\'nke yi Kitule}. \url{https://grambank.clld.org/languages/tula1252}
\item Glottolog reference for Johannes Lukas (1954/56), a grammar of Tula originally written by the missionary J. Hall. \url{https://glottolog.org/resource/reference/id/89075}
\item Wikipedia, ``Tula language''. It links Roger Blench's work on the phonology and noun morphology of Tula and his wordlist. \url{https://en.wikipedia.org/wiki/Tula_language}
\item Wikipedia, ``Tula--Waja languages''. \url{https://en.wikipedia.org/wiki/Tula%E2%80%93Waja_languages}
\item Adamawa Languages Project, Mainz: Tula-Waja group. It notes the complex noun class system. \url{https://www.blogs.uni-mainz.de/fb07-adamawa/tula-waja-group/}
\item Adamawa Grammar (AdaGram) page, LLACAN-CNRS. \url{https://llacan.cnrs.fr/AdaGram/dijim.html}
\end{itemize}

\section{Religious and Community Resources}
The search found no Tula Bible translation or community organisation. A Gombe State government page (\url{https://gombestate.gov.ng/?p=16209}) came up in the results, but its relevance to Tula was not established.

\section{Digital and Audio Resources}
No audio, video or software resources for Tula were found.

\section{Curriculum Development Context}
Tula is not known to be taught in schools, and no standardized teaching materials were found. A curriculum would have to start from the dictionaries and the Blench documentation. Local speakers and the dictionary authors would need to be consulted on the orthography and on dialect choice, since three dialects are reported. The speaker figure of about 30,000 comes from an old source.

\end{document}